\chapter{A Concrete Naming Certificate for $\BishopCauchyModulusCofinalityUp$}
\label{ch:concrete-instances-bishop-cauchy-modulus-cofinality-namecert}
\origin{human}

Following Bishop and Bridges, the $\BishopCauchyModulusCofinalityUp$ packet records the finite paired-modulus common-tail certificate used when two Bishop Cauchy schedules must be read through one later tail. It sits after the single-modulus schedule row of \autoref{ch:concrete-instances-cauchymodulus-namecert}, the dyadic tolerance ledger of \autoref{ch:concrete-instances-dyadicratcore-namecert}, the finite stream-window surface of \autoref{ch:concrete-instances-streamname-namecert}, the regular rational readback of \autoref{ch:concrete-instances-regseqrat-namecert}, the terminal real-name seal of \autoref{ch:concrete-instances-real-namecert}, and the cofinal-window route of \autoref{ch:concrete-instances-regular-cauchy-window-cofinality-namecert}. The packet is a paired finite tail ledger for L10 Real and Completion consumers; it is not quotient stream equality, a selected limit, countable choice, trichotomy, or ambient Real completeness.

\begin{definition}[BishopCauchyModulusCofinality carrier]
\label{def:bishop-cauchy-modulus-cofinality-carrier}
A $\BishopCauchyModulusCofinalityUp$ carrier is a finite $\BHist$ packet
$$
M=(\epsilon,D,S_0,S_1,R_0,R_1,\mu_0,\mu_1,W,T,E,H,C,P,N)
$$
with the following displayed rows:
\begin{enumerate}
\item $\epsilon$ is the requested tolerance row.
\item $D$ is the $\DyadicRatCoreUp$ comparison ledger for $\epsilon$ and the two tail budgets.
\item $S_0$ and $S_1$ are the two $\StreamNameUp$ finite-window rows being compared.
\item $R_0$ and $R_1$ are the matching $\RegSeqRatUp$ readback rows for those windows.
\item $\mu_0$ and $\mu_1$ are the two $\CauchyModulusUp$ schedule rows.
\item $W$ is the $\RegularCauchyWindowCofinalityUp$ row that displays a later finite window cofinal for both schedules at tolerance $\epsilon$.
\item $T$ is the common-tail certificate, stating that every read through $W$ factors through both $\mu_0$ and $\mu_1$ after the displayed dyadic comparison.
\item $E$ is the $\RealUp$ seal row reached only after $D,S_0,S_1,R_0,R_1,\mu_0,\mu_1,W,T$ have been displayed.
\item $H$ records componentwise $\hsame$ transport for tolerances, dyadic comparisons, stream windows, regular readbacks, Cauchy moduli, cofinal windows, common-tail certificates, Real seals, provenance, and names.
\item $C$ records displayed $\Cont$ replay from the two modulus schedules to $W,T,E$ through the same finite rows.
\item $P$ is the $\Pkg$ provenance over $\BishopCauchyModulusCofinalityUp$, $\CauchyModulusUp$, $\DyadicRatCoreUp$, $\StreamNameUp$, $\RegSeqRatUp$, $\RegularCauchyWindowCofinalityUp$, $\RealUp$, $\BHist$, $\hsame$, $\Cont$, and $\NameCert$.
\item $N$ is the local $\NameCert_{\BishopCauchyModulusCofinalityUp}$ row.
\end{enumerate}
The classifier compares accepted packets only through these rows. It has no coordinate for quotient streams, selected limits, countable choice, trichotomy, or ambient Real completeness.
\end{definition}

\begin{theorem}[BishopCauchyModulusCofinality NameCert obligations]
\label{thm:bishop-cauchy-modulus-cofinality-namecert-obligations}
For every accepted $\BishopCauchyModulusCofinalityUp$ carrier $M=(\epsilon,D,S_0,S_1,R_0,R_1,\mu_0,\mu_1,W,T,E,H,C,P,N)$, the local $\NameCert_{\BishopCauchyModulusCofinalityUp}$ surface consists exactly of tolerance admission through $\epsilon$; dyadic comparison through $D$; two finite $\StreamNameUp$ windows through $S_0,S_1$; two $\RegSeqRatUp$ readbacks through $R_0,R_1$; two $\CauchyModulusUp$ schedules through $\mu_0,\mu_1$; cofinal-window exposure through $W$; common-tail exposure through $T$; $\RealUp$ sealing through $E$; componentwise transport through $H$; displayed replay through $C$; provenance through $P$; and local naming through $N$.
\end{theorem}
\begin{proof}
Project the carrier rows from \autoref{def:bishop-cauchy-modulus-cofinality-carrier}. The accepted packet displays one tolerance, one dyadic comparison ledger, two stream windows, two regular rational readbacks, two Cauchy modulus schedules, one cofinal window, one common-tail certificate, one Real seal, and the structural transport, replay, provenance, and naming rows.

The paired nature of the certificate is carried only by $\mu_0,\mu_1,W,T$. The cofinal-window row $W$ is read after the two schedules and the dyadic ledger, and $T$ records that the later finite tail factors through both schedules.

The structural rows preserve the same finite packet under $\hsame$, $\Cont$, $\Pkg$, and $\NameCert$ replay. Since no row carries quotient stream equality, a selected limit, countable choice, trichotomy, or ambient Real completeness, the listed rows exhaust the NameCert obligations.
\end{proof}

\begin{theorem}[BishopCauchyModulusCofinality common tail]
\label{thm:bishop-cauchy-modulus-cofinality-common-tail}
If an accepted $\BishopCauchyModulusCofinalityUp$ packet presents two Cauchy modulus schedules $\mu_0,\mu_1$ over the displayed $\StreamNameUp$ and $\RegSeqRatUp$ rows, then every common-tail read at tolerance $\epsilon$ factors through the finite route
$$
D,S_0,S_1,R_0,R_1,\mu_0,\mu_1,W,T.
$$
The common tail is the displayed cofinal window $W$ together with certificate $T$, not a quotient-stream identification, selected limit, countable-choice selection, trichotomy branch, or Real-completeness argument.
\end{theorem}
\begin{proof}
Begin with the dyadic tolerance row $D$. It fixes the comparison budget for the two schedules and prevents the common-tail read from changing tolerance outside the displayed packet.

Next read the two stream and regular-rational sides. The rows $S_0,S_1$ expose the finite windows, and $R_0,R_1$ attach those windows to the regular rational readbacks that the two modulus schedules govern.

Finally apply the cofinal-window row $W$ and the common-tail certificate $T$. These rows display the later finite tail that is after both schedules at the requested tolerance, so a consumer receives a paired finite tail certificate without invoking quotient streams, selected limits, countable choice, trichotomy, or Real completeness.
\end{proof}

\begin{theorem}[BishopCauchyModulusCofinality Real-seal handoff]
\label{thm:bishop-cauchy-modulus-cofinality-real-seal-handoff}
Every $\RealUp$-facing read of an accepted $\BishopCauchyModulusCofinalityUp$ packet reaches the seal $E$ only after the dyadic ledger, both stream windows, both regular rational readbacks, both Cauchy modulus schedules, the $\RegularCauchyWindowCofinalityUp$ row, and the common-tail certificate have been displayed. The handoff exports a finite paired-modulus tail certificate for L10 Real and Completion routes, and exports no quotient stream equality, selected limit, countable choice, trichotomy, or ambient Real completeness.
\end{theorem}
\begin{proof}
Project the finite route before the seal. The rows $D,S_0,S_1,R_0,R_1,\mu_0,\mu_1,W,T$ are all upstream of $E$, so the Real-facing consumer cannot see the seal without also seeing the paired-modulus common-tail data.

Transport and replay preserve that order. The rows $H$ and $C$ can carry the displayed packet through equivalent naming and continuation reads, but they do not add another stream quotient, hidden tail selector, or completeness coordinate.

Thus the terminal $\RealUp$ read is only a seal for the finite paired-tail route. It remains choice-free and quotient-free, and it does not decide trichotomy or claim a selected completion limit.
\end{proof}

\falsifiablePrediction{An accepted $\BishopCauchyModulusCofinalityUp$ consumer that reaches a Real-facing common tail without displaying the dyadic ledger, both StreamName windows, both RegSeqRat readbacks, both CauchyModulus schedules, the RegularCauchyWindowCofinality row, the common-tail certificate, and the structural rows refutes the carrier surface.}

\independenceWitness{cauchymodulus, dyadicratcore, streamname, regseqrat, regular-cauchy-window-cofinality, real: no listed sibling supplies the whole paired packet $M=(\epsilon,D,S_0,S_1,R_0,R_1,\mu_0,\mu_1,W,T,E,H,C,P,N)$, because each sibling contributes only one schedule, dyadic ledger, stream window, regular readback, cofinal-window row, or terminal seal while this chapter binds two modulus schedules into one common-tail certificate.}

\closureat{BishopCauchyModulusCofinalityUp}{seedStr}

\begin{closurestatus}{\BishopCauchyModulusCofinalityUp}
  \constructivestory{A $\BishopCauchyModulusCofinalityUp$ packet is generated from one dyadic tolerance ledger, two StreamName windows, two RegSeqRat readbacks, two CauchyModulus schedules, one RegularCauchyWindowCofinality row, a common-tail certificate, a RealUp seal, componentwise hsame transport, displayed Cont replay, Pkg provenance, and a local NameCert row.}
  \theoryclosure{\seedClosure}
  \scopeclosed{\autoref{def:bishop-cauchy-modulus-cofinality-carrier}, \autoref{thm:bishop-cauchy-modulus-cofinality-namecert-obligations}, \autoref{thm:bishop-cauchy-modulus-cofinality-common-tail}, and \autoref{thm:bishop-cauchy-modulus-cofinality-real-seal-handoff} seed the finite paired-modulus common-tail route for L10 Real and Completion consumers.}
  \formalstatus{\unformalizedV}
  \bridgestatus{none}
  \notclaimed{This chapter does not close quotient stream equality, selected limits, countable choice, trichotomy, ambient Real completeness, Lean verification, or bridge checking.}
  \upgradepath{Obligation closure requires checked carrier admission, paired schedule exposure, cofinal-window exactness, common-tail stability, Real-seal nonescape, transport, replay, provenance, and local naming for BEDC.Derived.BishopCauchyModulusCofinalityUp.}
\end{closurestatus}
